\chapter{完整源码清单：Cpp 文本分析工作簿}

本附录把 \filepath{books/cpp-zero-to-advanced/labs/cpp_practice_workbook} 的完整主线工程写入书中。阅读时按“构建入口、公共接口、库实现、CLI、测试、benchmark”的顺序走。每个文件前面的说明告诉你它解决什么问题、不要让它承担什么责任。

\section{一、CMakeLists.txt}

CMake 文件把库、命令行程序、测试和 benchmark 分成不同 target。它是工程边界的第一层：库不依赖 CLI，测试链接库，benchmark 只测核心路径。新增源码时先检查它应该属于哪个 target。

\lstinputlisting[language=bash,caption={cpp\_practice\_workbook/CMakeLists.txt}]{../../labs/cpp_practice_workbook/CMakeLists.txt}

\section{二、README}

README 是读者复现实验的最短路径。它列出构建、测试和运行命令；如果命令和书中正文不一致，应以可执行工程为准并同步修正文档。

\lstinputlisting[language=bash,caption={cpp\_practice\_workbook/README.md}]{../../labs/cpp_practice_workbook/README.md}

\section{三、公共头文件}

\filepath{practice_workbook.hpp} 定义本工程的公共数据合同：配置、诊断、词频结果、分析器和缓存。读者应先理解这些类型，再进入实现文件；如果头文件暴露过多实现细节，后续章节的抽象边界就会变差。

\lstinputlisting[language=C++,caption={include/cppbook/practice\_workbook.hpp}]{../../labs/cpp_practice_workbook/include/cppbook/practice_workbook.hpp}

\section{四、库实现}

实现文件包含配置解析、文本扫描、词频排序、LRU 缓存和文件读取。读代码时按函数责任分层：字符串清洗不做文件 I/O，解析器不退出进程，缓存只维护不变量，文件读取只负责 RAII 资源边界。

\lstinputlisting[language=C++,caption={src/practice\_workbook.cpp}]{../../labs/cpp_practice_workbook/src/practice_workbook.cpp}

\section{五、命令行入口}

\filepath{src/main.cpp} 应保持很薄：解析命令行、读取文件、调用库、打印结果。它不能复制 \code{TextAnalyzer} 或配置解析逻辑，否则库测试通过也不能保证 CLI 行为正确。

\lstinputlisting[language=C++,caption={src/main.cpp}]{../../labs/cpp_practice_workbook/src/main.cpp}

\section{六、完整测试}

测试文件覆盖配置成功、配置失败、文本分析、数字保留、缓存淘汰、文件读取和 limit 参数。读者应把每个测试名当成行为合同：以后重构时，测试名不应该失去可读性。

\lstinputlisting[language=C++,caption={tests/practice\_workbook\_tests.cpp}]{../../labs/cpp_practice_workbook/tests/practice_workbook_tests.cpp}

\section{七、Benchmark}

benchmark 只回答一个问题：核心文本分析路径随输入规模怎样变化。它不把文件读取、打印和配置解析混进计时区域；否则性能结论会失去解释力。

\lstinputlisting[language=C++,caption={bench/practice\_workbook\_bench.cpp}]{../../labs/cpp_practice_workbook/bench/practice_workbook_bench.cpp}
